\documentclass[book.tex]{subfiles}
\begin{document}

\chapter{Euclidean Networks as an Equivariant Framework}
\label{ch:dft_e3nn}

%\url{https://scholar.google.com/citations?user=iF01q24AAAAJ&hl=en&oi=sra}
\label{chap:e3nn}
\noindent\rule{11cm}{0.4pt} \\
\textbf{Prerequisites:} \autoref{chap:dft}, \autoref{chap:equivariant_networks} \\
\textbf{Difficulty Level:} ** \\
\noindent\rule{11cm}{0.4pt}
\vspace{5mm}

%"""
%Equivariant neural networks (ENNs) are graph neural networks embedded in R3 and are well suited for predicting molecular properties. The ENN library e3nn has customizable convolutions, which can be designed to depend only on distances between points, or also on angular features, making them rotationally invariant, or equivariant, respectively. This paper studies the practical value of including angular dependencies for molecular property prediction directly via an ablation study with e3nn and the QM9 data set. We find that, for fixed network depth and parameter count, adding angular features decreased test error by an average of 23\%. Meanwhile, increasing network depth decreased test error by only 4\% on average, implying that rotationally equivariant layers are comparatively parameter efficient. We present an explanation of the accuracy improvement on the dipole moment, the target which benefited most from the introduction of angular features.
%"""
%References 

Euclidean networks bridge the several different types of equivariant convolutions we have seen (tensor field networks, $SE(3)$ transformers, steerable convolutional networks) into one common framework.

\section{Mathematical Framework}

Suppose that atomic positions are given by
\begin{align}
r \in \mathbb{R}^{3N}
\end{align}
where $N$ is the number of atoms. In order to define more complex equivariant architectures, we need the ability to manipulate a few basic ingredients. First, we need efficient implementations of the spherical harmonics and irreducible representations of common groups. We also need precomputed implementations of the Clebsch-Gordon coefficients. Given these core building blocks, it becomes possible to build a large range of equivariant networks.
%Equivariant convolutions \cite{miller2020relevance} can be used to model complex molecular potentials.

%\begin{align}
%    f_i' &= \frac{1}{\sqrt{z}} \sum_{j \in \partial(i)} f_i \otimes(h(\|x_{ij}\|)) Y(x_{ij}/\|x_{ij}\|)
%\end{align}

%Here $f_i$ are the input nodes, $f_i'$ are the output nodes, $z$ is the average degree of the nodes, $\partial(i)$ is the set of neighbors of node $i$, $x_{ij}$ is the relative vector, $h$ is a multi-layer perceptron, $Y$ is a spherical harmonic, and $x \times(w) y$ is a tensor product of $x$ with $y$ parameterized by weights $w$  

\end{document}